\documentclass[10pt,a4paper]{amsart}
\usepackage[margin=19mm]{geometry}
\usepackage{amsmath,amssymb,mathtools}
\usepackage[scr=rsfs,cal=euler]{mathalfa}
\usepackage{xfrac}
\usepackage[unicode,hidelinks,bookmarks=false]{hyperref}
\newcommand{\ie}{i.e.\ }
\newcommand{\cf}{cf.\ }
\newcommand{\resp}{respectively\ }
\newcommand{\Subsection}{\subsection}
\providecommand{\nobreakdash}{\nobreak\hbox{-}}
\setlength{\emergencystretch}{2em}
\multlinegap0pt
\usepackage{amssymb}
\usepackage{tikz}
\newtheorem{definition}{Definition}
\newtheorem{proposition}{Proposition}
\newcommand{\Aut}{\mathrm{Aut}}
\newcommand{\C}{{\mathrm{C}}}
\newcommand\DD{{\mathscr{D}}}
\title{Quasiprimitive and bi-quasiprimitive highly-arc-transitive digraphs and finite simple groups}
\author{Lei Chen, Cheryl Praeger}
\date{Source: arXiv:2512.17244v1 (CC BY 4.0)}
\begin{document}
\maketitle

\noindent\emph{Source and licence.} Quasiprimitive and bi-quasiprimitive highly-arc-transitive digraphs and finite simple groups. Version arXiv:2512.17244v1 (CC BY 4.0), distributed by arXiv under the Creative Commons Attribution 4.0 International licence, http://creativecommons.org/licenses/by/4.0/, as recorded in the arXiv licence metadata for this exact version and shown on the abstract page https://arxiv.org/abs/2512.17244v1. Full licence notice: \texttt{licenses/arxiv-2512.17244-CC-BY-4.0.txt}.\par\smallskip

\noindent\textit{Editorial scope.} Counted unit: the article's proposition p:quot2 (source0.tex lines 354--361). The article's complete original proof of it is delivered with it (source0.tex lines 363--369, one \texttt{proof} environment of 7 lines). No mathematics is written, repaired or re-derived: the statement and the proof are the article's own lines, in the article's own notation, and no retained line is substituted or re-flowed.  Retained uncounted as the source-local context the proof needs: lines 319--325; lines 336--342.  Nothing was omitted from any retained span, so the package carries no omission record.  No retained line contains a citation, so the wrapper carries no bibliography and no bibliography entry.  No retained line contains a figure environment, an image-inclusion command or drawing code, so no image is delivered and the package declares no asset dependency.  Editorial formatting only, in addition: an AMS \texttt{amsart} preamble that restates the article's own declarations and macros used by the retained lines, namely the environments \texttt{definition}, \texttt{proposition} on separate counters, together with the standard packages those lines need.  The retained environments print in this document as Definition 1, Proposition 1 in order of appearance, whereas the article prints the same items with the numbering of its own full document.  The complete native source of the article as distributed in the arXiv e-print for this exact version is bundled unchanged as \texttt{sources/191302/source0.tex}.
\begin{definition}\label{d:nquot}
{\rm 
    Let $V_N$ denote that set of $N$-orbits in $V$ and define the \emph{quotient} \(\Gamma_N\) to have vertex set $V_N$ and arc set $A_N\subseteq V_N\times V_N$ consisting of all pairs $(U,U')\in V_N\times V_N$ such that,  for some $u\in U, u'\in U'$, the pair $(u,u')\in A$.  

    In particular suppose that there exists a vertex-transitive subgroup $G\leq \Aut(\Gamma)$ containing $N$. If $N\unlhd G$, then $\Gamma_N$ is called the \emph{$G$-normal quotient of $\Gamma$ relative to $N$}. Alternatively, if $N \unlhd M\unlhd G$, $N$ is not normal in $G$, and $M$ is transitive on $V$, then $\Gamma_N$ is called a \emph{$G$-subnormal quotient of $\Gamma$ relative to $N$}.
    }
\end{definition}

\begin{proposition}\label{p:quot}
   Let \(\Gamma=(V,A)\) be a finite connected $G$-vertex-transitive  \((G,s)\)-arc-transitive digraph with \(s\geqslant1\) and $G\leq\Aut(\Gamma)$. Let $N\leq G$ with $|V_N|\geq3$ and $\Gamma_N$ as in Definition~\ref{d:nquot}. Then $\Gamma_N$ is connected, and moreover the following hold.
   \begin{enumerate}
       \item[(a)] If $N\unlhd G$ and  \(s\geqslant2\), then $\Gamma_N$ is a digraph and is $G/N$-vertex-transitive and $(G/N,s)$-arc-transitive. 
       \item[(b)] If $N\unlhd M\unlhd G$ and  \(s\geqslant3\), with $N$ not normal in $G$ and $M$ vertex-transitive on $\Gamma$, then the $G$-subnormal quotient $\Gamma_N$ is a digraph and is $N_G(N)/N$-vertex-transitive and  $(N_G(N)/N,s-1)$-arc-transitive. 
   \end{enumerate} 
\end{proposition}

\begin{proposition}\label{p:quot2}
   Let \(\Gamma=(V,A)\) be a finite connected $G$-vertex-transitive \((G,s)\)-arc-transitive digraph with \(s\geqslant2\) and $G\leq\Aut(\Gamma)$. Let $N\unlhd G$ with $|V_N|\geq3$. Then the following hold.
   \begin{enumerate}
       \item[(a)] If  $N\leqslant M\unlhd G$ with $|V_M|\geqslant3$, then $M/N\unlhd G/N$ and $\Gamma_M\cong (\Gamma_N)_{M/N}$.
       \item[(b)] If $r\geqslant3$ and no $G$-normal quotient of $\Gamma$ is isomorphic to $\C^{\to}_r$, then  no $G/N$-normal quotient of $\Gamma_N$ is isomorphic to $\C^{\to}_r$.
       \item[(c)] There exists $M\unlhd G$ such that $N\leqslant M$ and $|V_M|\geqslant3$, and the $G$-normal quotient $\Gamma_M$ is a $(G/M,s)$-arc-transitive digraph with $G/M$ quasiprimitive or bi-quasiprimitive on $V_M$. Moreover, if $(\Gamma, G, s)\in\DD$ then also $(\Gamma_M, G/M, s)\in\DD$.
   \end{enumerate}
\end{proposition}

\begin{proof}
    (a) Suppose that $N\leqslant M\unlhd G$ with $|V_M|\geqslant3$. Then $M/N\unlhd G/N$ and each $M$-orbit $u^M\in V_M$ is a union of all the $N$-orbits of the form $(u^N)^x=u^{Nx}$ for $x\in M$, and it follows that the map $\phi:V_M\to (V_N)_{M/N}$ sending $u^M$ to the $M/N$-orbit in $V_N$ containing $u^N$ is a bijection. Also by Definition~\ref{d:nquot}, $\phi$ induces a bijection from the arcs of $\Gamma_M$ to the arcs of $(\Gamma_N)_{M/N}$, giving a graph isomorphism $\Gamma_M\to  (\Gamma_N)_{M/N}$.

    (b) If $(\Gamma_N)_{\overline{M}}\cong \C^{\to}_r$ for some normal subgroup $\overline{M}=M/N$ of $G/N$, then $N\leqslant M\unlhd G$ and by part (a), the $G$-normal quotient $\Gamma_M$ is isomorphic to $\C^{\to}_r$. Thus part (b) is proved.

    (c) Choose a subgroup $M$ that is maximal by inclusion such that $N\leqslant M$ and $|V_M|\geqslant 3$. Such a subgroup exists because $N$ has these two properties and $G$ is finite. By Proposition~\ref{p:quot}(a), $\Gamma_M$ is a $(G/M,s)$-arc-transitive digraph. Further, each nontrivial normal subgroup of $G/M$ is of the form $L/M$ for some $L\unlhd G$ properly containing $M$. It follows from the maximality of $M$ that $|V_L|\leq 2$. Thus $L/M$ has at most two orbits in $V_M$ and so $G/M$ is either quasiprimitive or bi-quasiprimitive on $V_M$.  Finallly,  if $(\Gamma, G, s)\in\DD$ then  no $G$-normal quotient of $\Gamma$ is isomorphic to $\C^{\to}_r$ for any $r\geq3$, and by part (b),  no $G/M$-normal quotient of $\Gamma_M$ is isomorphic to $\C^{\to}_r$ for any $r\geq3$. Hence $(\Gamma_M, G/M, s)\in\DD$.
\end{proof}

\end{document}
